% Instructor-authored original practice bank. This is an includable LaTeX fragment.
% Each source-map comment records only the disclosed form location and an abstract skill.

\section*{2025 Logical Reasoning: Original Analogues}

\begin{center}
\fbox{\parbox{0.92\linewidth}{\textbf{Instructor-authored original material.}
The names LSAT and LSAC are used only to identify general skill inspiration.
Every scenario, claim, question, and answer choice below is newly written.
This material is not written, approved, or endorsed by LSAC.}}
\end{center}

\noindent This is a public practice bank, not a live graded form. Choose the best answer to each question.

\subsection*{Practice Set 1}

\begin{enumerate}

% Source map: 2025 LR Section 1, Question 01 — principle support
\item A reviewer says that a museum should not reject a sculpture merely because its message is unpopular. A museum that promises to present serious contemporary work should judge submissions by artistic merit, provided that displaying them would not create a safety hazard.

Which principle most strongly supports the reviewer's reasoning?
\begin{enumerate}
\item[(A)] Museums should exhibit every work submitted by a professional artist.
\item[(B)] An institution should apply its stated selection criteria even when doing so produces controversy.
\item[(C)] Public institutions should avoid taking positions on disputed questions.
\item[(D)] Safety concerns are never relevant to artistic decisions.
\item[(E)] A museum's most popular exhibits are generally its most valuable ones.
\end{enumerate}

% Source map: 2025 LR Section 1, Question 02 — necessary assumption
\item A warehouse declined to replace its experienced night crew with cheaper temporary workers. Managers are under orders to cut costs wherever doing so would not risk delaying shipments. Therefore, the managers must believe that making the replacement could delay shipments.

The explanation requires assuming which statement?
\begin{enumerate}
\item[(A)] Temporary workers are never as productive as experienced workers.
\item[(B)] Night shipments are more important than daytime shipments.
\item[(C)] The warehouse has already missed several shipments.
\item[(D)] The managers generally follow the order to cut costs when its stated condition is met.
\item[(E)] The experienced crew requested higher wages.
\end{enumerate}

% Source map: 2025 LR Section 1, Question 03 — point at issue
\item Mina says, ``The seminar should include the difficult article because it introduces the central vocabulary.'' Rowan replies, ``Students can learn that vocabulary from the handbook, so the difficult article is unnecessary.''

Mina and Rowan disagree about whether
\begin{enumerate}
\item[(A)] the handbook defines the central vocabulary accurately.
\item[(B)] students should learn any vocabulary before the seminar.
\item[(C)] including the difficult article is necessary for the seminar.
\item[(D)] the article is more difficult than the handbook.
\item[(E)] seminars should include assigned reading.
\end{enumerate}

% Source map: 2025 LR Section 1, Question 04 — flaw
\item A maker of reusable bottles claims its bottles never leak because no customer has proved that a bottle leaked when its lid was correctly tightened. Therefore, every report of leakage must involve a lid that was not correctly tightened.

The reasoning is flawed because it
\begin{enumerate}
\item[(A)] treats a necessary condition for leakage as sufficient.
\item[(B)] treats the absence of conclusive evidence for a claim as conclusive evidence against it.
\item[(C)] generalizes from a sample known to be unrepresentative.
\item[(D)] attacks the motives of dissatisfied customers.
\item[(E)] assumes that reusable bottles are preferable to disposable ones.
\end{enumerate}

% Source map: 2025 LR Section 1, Question 05 — sufficient assumption
\item Every mural on the east wall was painted after 1990. The blue mural is on the east wall. Therefore, the blue mural is newer than the red mural.

The conclusion follows logically if which assumption is added?
\begin{enumerate}
\item[(A)] The red mural was painted no later than 1990.
\item[(B)] The blue mural was restored after it was painted.
\item[(C)] Some murals not on the east wall were painted after 1990.
\item[(D)] The red mural is not on the east wall.
\item[(E)] The east wall contains at least two murals.
\end{enumerate}

% Source map: 2025 LR Section 1, Question 06 — weaken
\item An archivist argues that the small holes along the edges of certain paper folders serve no current function because the folders are now stored flat rather than hung from hooks. Thus, the holes must merely be a remnant of an older filing system.

Which fact most weakens the argument?
\begin{enumerate}
\item[(A)] Some folders were manufactured before the archive opened.
\item[(B)] The archive uses more than one folder size.
\item[(C)] Similar holes appear in folders at other archives.
\item[(D)] The archive's digitization machine uses the holes to align documents during scanning.
\item[(E)] Hooks remain in one unused storage room.
\end{enumerate}

% Source map: 2025 LR Section 1, Question 07 — resolve discrepancy
\item A ferry operated fewer hours in May than in April, yet it carried substantially more passengers in May. Its fare and passenger capacity did not change.

Which fact best resolves the apparent discrepancy?
\begin{enumerate}
\item[(A)] The ferry carried less freight in May.
\item[(B)] Its crew worked the same shifts in both months.
\item[(C)] A commuter-rail station opened beside the ferry terminal in May, bringing many new riders.
\item[(D)] Fuel cost the ferry operator more in May.
\item[(E)] The ferry's timetable used a new typeface in May.
\end{enumerate}

% Source map: 2025 LR Section 1, Question 08 — role in argument
\item Restoring the marsh will reduce flooding, since intact marshes absorb storm water. Moreover, the restoration will create bird habitat. The city should therefore fund the project this year.

What role does the claim that intact marshes absorb storm water play?
\begin{enumerate}
\item[(A)] It is the argument's main conclusion.
\item[(B)] It supports an intermediate claim that helps justify the main conclusion.
\item[(C)] It states an objection that the argument rejects.
\item[(D)] It provides background unrelated to the recommendation.
\item[(E)] It follows from the recommendation rather than supporting it.
\end{enumerate}

% Source map: 2025 LR Section 1, Question 09 — weaken
\item A city installed bright lane markers at five dangerous intersections. Collisions at those intersections declined the next month. Officials conclude that the lane markers caused the decline.

Which fact most weakens the conclusion?
\begin{enumerate}
\item[(A)] The markers cost less than officials expected.
\item[(B)] Drivers generally say that bright markings are easy to see.
\item[(C)] Collision rates vary somewhat from month to month.
\item[(D)] Two nearby cities use similar markers.
\item[(E)] During the same month, construction closed one approach to each intersection and sharply reduced traffic.
\end{enumerate}

% Source map: 2025 LR Section 1, Question 10 — parallel flaw
\item Since every cedar in the park is a tree and some trees in the park are diseased, some cedars in the park must be diseased.

Which argument contains the same flaw?
\begin{enumerate}
\item[(A)] Every violin is an instrument, and some instruments in the shop are damaged; therefore, some violins in the shop are damaged.
\item[(B)] Some buses are electric, and every electric vehicle is quiet; therefore, some buses are quiet.
\item[(C)] No ceramic cups are flexible, and this cup is ceramic; therefore, this cup is not flexible.
\item[(D)] Every comet orbits the Sun, and object K is a comet; therefore, object K orbits the Sun.
\item[(E)] Some editors are poets, and some poets are teachers; therefore, some editors may be teachers.
\end{enumerate}

% Source map: 2025 LR Section 1, Question 11 — resolve discrepancy
\item A wetland survey recorded fewer frogs this spring than last spring, even though rainfall was greater and the wetland's total frog population increased.

Which fact best resolves the apparent discrepancy?
\begin{enumerate}
\item[(A)] Frogs require water in order to reproduce.
\item[(B)] The survey team used the same notebooks both years.
\item[(C)] Several frog species live in the wetland.
\item[(D)] High water created many new pools outside the fixed area covered by the survey.
\item[(E)] The wetland also supports salamanders.
\end{enumerate}

% Source map: 2025 LR Section 1, Question 12 — strengthen
\item After a ceramics studio installed a new ventilation system, reports of headaches among its evening students decreased. A researcher concludes that the ventilation change caused the decrease.

Which fact most strengthens the conclusion?
\begin{enumerate}
\item[(A)] Some students prefer working with porcelain.
\item[(B)] Comparable nearby studios that did not change ventilation showed no similar decrease during the same period.
\item[(C)] The studio replaced its kilns several years earlier.
\item[(D)] Evening classes are shorter than weekend workshops.
\item[(E)] Students report symptoms electronically.
\end{enumerate}

% Source map: 2025 LR Section 1, Question 13 — main conclusion
\item A diagram cannot contain every detail of the system it represents. Adding every detail would make the diagram as complicated as the system and therefore useless as a guide. Thus, a useful diagram must omit some true information, even though those omissions can occasionally mislead.

What is the main conclusion?
\begin{enumerate}
\item[(A)] Diagrams occasionally mislead.
\item[(B)] A system can be more complicated than its diagram.
\item[(C)] A useful diagram must omit some true information.
\item[(D)] No diagram can accurately represent a system.
\item[(E)] Guides should contain every available detail.
\end{enumerate}

% Source map: 2025 LR Section 1, Question 14 — principle identification
\item A committee decides to publish the evidence behind its recommendation before asking the public to accept that recommendation.

Which principle is most clearly illustrated?
\begin{enumerate}
\item[(A)] People asked to accept a decision should be able to inspect the reasons offered for it.
\item[(B)] Committees should never revise a recommendation after publishing it.
\item[(C)] Public approval is the only measure of a recommendation's quality.
\item[(D)] Evidence is useful only when everyone interprets it alike.
\item[(E)] A recommendation should be delayed until no one objects.
\end{enumerate}

% Source map: 2025 LR Section 1, Question 15 — point at issue
\item Sato says, ``Requiring identical signs on every shop would make the street easier to navigate.'' Bell replies, ``That benefit does not justify erasing the distinctive designs that make the street inviting.''

Sato and Bell most clearly disagree about whether
\begin{enumerate}
\item[(A)] distinctive signs can help customers navigate.
\item[(B)] every existing sign is attractive.
\item[(C)] the street is currently difficult to navigate.
\item[(D)] easier navigation justifies requiring identical signs.
\item[(E)] shops should display signs.
\end{enumerate}

% Source map: 2025 LR Section 1, Question 16 — role in argument
\item Some claim that a machine that predicts our choices would show that deliberation is pointless. But a prediction can incorporate the fact that people deliberate, just as a weather forecast incorporates atmospheric motion. The claim that prediction makes deliberation pointless is therefore mistaken.

What role does the comparison with a weather forecast play?
\begin{enumerate}
\item[(A)] It states the conclusion the argument seeks to establish.
\item[(B)] It illustrates how prediction can depend on, rather than replace, the process predicted.
\item[(C)] It concedes that human choices are identical to weather events.
\item[(D)] It introduces evidence that weather cannot be predicted.
\item[(E)] It defines what counts as deliberation.
\end{enumerate}

% Source map: 2025 LR Section 1, Question 17 — sufficient assumption
\item If the supplier closes, the workshop will miss its deadline unless it has enough stored oak. The supplier will close. Therefore, the workshop will miss its deadline.

Which assumption makes the conclusion follow logically?
\begin{enumerate}
\item[(A)] Oak is the workshop's most expensive material.
\item[(B)] The supplier has closed before.
\item[(C)] The workshop does not have enough stored oak.
\item[(D)] The deadline cannot be extended.
\item[(E)] The workshop uses no other supplier.
\end{enumerate}

% Source map: 2025 LR Section 1, Question 18 — parallel reasoning
\item Whenever the library hosts an evening lecture, it keeps the west door open. The west door is not open tonight. Therefore, the library is not hosting an evening lecture tonight.

Which argument uses the same pattern of reasoning?
\begin{enumerate}
\item[(A)] If the kiln is hot, the studio is open. The studio is open, so the kiln is hot.
\item[(B)] If the pond freezes, skating is allowed. The pond froze, so skating is allowed.
\item[(C)] If the bell rings, class ends. Class did not end, but the bell may have rung.
\item[(D)] If the archive closes, researchers leave. Researchers left, so the archive closed.
\item[(E)] If the greenhouse is being watered, its pump is running. The pump is not running, so the greenhouse is not being watered.
\end{enumerate}

% Source map: 2025 LR Section 1, Question 19 — most strongly supported
\item Prototype batteries generate more heat than conventional batteries when discharged rapidly. In cold-weather trials, high-drain sensors using the prototypes stop sooner than comparable sensors in mild conditions, while low-drain sensors show little difference.

Which statement is most strongly supported?
\begin{enumerate}
\item[(A)] Cold conditions can constrain prototype batteries more during rapid discharge than during slow discharge.
\item[(B)] Conventional batteries generate no heat.
\item[(C)] Every sensor that stops early uses a prototype battery.
\item[(D)] Warming any battery always increases its total lifetime.
\item[(E)] Low-drain sensors never operate in cold weather.
\end{enumerate}

% Source map: 2025 LR Section 1, Question 20 — main conclusion
\item Critics say that public murals are valuable only when they last for decades. Yet temporary murals can invite experimentation, and their replacement can bring new artists into public view. Therefore, a mural's short life does not by itself make the project artistically wasteful.

What is the main conclusion?
\begin{enumerate}
\item[(A)] All temporary murals are artistically successful.
\item[(B)] New artists prefer temporary work.
\item[(C)] Long-lasting murals discourage experimentation.
\item[(D)] A short lifespan alone does not make a mural project artistically wasteful.
\item[(E)] Public art should be replaced frequently.
\end{enumerate}

% Source map: 2025 LR Section 1, Question 21 — flaw
\item An ecologist notes that blue streetlights are common in districts where moth counts have fallen. The ecologist concludes that blue streetlights must be a major cause of the decline.

The argument is most vulnerable because it
\begin{enumerate}
\item[(A)] assumes without evidence that moths can detect light.
\item[(B)] rejects a claim merely because it is unpopular.
\item[(C)] infers causation from an association without ruling out other explanations.
\item[(D)] treats a sufficient condition as necessary.
\item[(E)] overlooks the fact that some streets lack sidewalks.
\end{enumerate}

% Source map: 2025 LR Section 1, Question 22 — must be true
\item Every accredited sanctuary employs a veterinarian. No organization that trades animals for profit is accredited. Meadow House is accredited, and River Yard trades animals for profit.

Which statement must be true?
\begin{enumerate}
\item[(A)] River Yard employs no veterinarian.
\item[(B)] Meadow House employs a veterinarian.
\item[(C)] Meadow House never transfers an animal.
\item[(D)] Every organization with a veterinarian is accredited.
\item[(E)] River Yard is the only unaccredited organization.
\end{enumerate}

% Source map: 2025 LR Section 1, Question 23 — flaw
\item A principal observes that every student in the robotics club takes algebra. Since Lee takes algebra, the principal concludes that Lee is in the robotics club.

The reasoning is flawed because it
\begin{enumerate}
\item[(A)] treats a condition required of club members as enough to establish membership.
\item[(B)] assumes that no club member takes geometry.
\item[(C)] draws a general conclusion from one case.
\item[(D)] relies on a term used in two different senses.
\item[(E)] confuses a group with one of its members.
\end{enumerate}

% Source map: 2025 LR Section 1, Question 24 — principle application
\item Principle: Borrowed equipment should be returned at least as functional as it was when received, unless the lender agrees otherwise.

Which action best conforms to the principle?
\begin{enumerate}
\item[(A)] A hiker returns a borrowed compass with a cracked face because it still points north.
\item[(B)] A musician returns a borrowed amplifier missing a knob after offering to replace it next year.
\item[(C)] A carpenter keeps a borrowed saw because the lender has not requested it back.
\item[(D)] A photographer returns a borrowed tripod with a loose leg after cleaning its case.
\item[(E)] A filmmaker replaces a frayed strap on a borrowed camera before returning it, without reducing any other function.
\end{enumerate}

% Source map: 2025 LR Section 1, Question 25 — flaw
\item Most students sampled from three evening courses commute by bus. A dean concludes that every student enrolled in an evening course commutes by bus.

The reasoning is most vulnerable because it
\begin{enumerate}
\item[(A)] assumes that buses operate only in the evening.
\item[(B)] relies on evidence gathered after the courses began.
\item[(C)] presumes that every sampled student answered honestly.
\item[(D)] infers a universal claim about all evening students from a majority in a limited sample.
\item[(E)] mistakes a claim about transportation for a claim about scheduling.
\end{enumerate}

\end{enumerate}

\subsection*{Practice Set 2}

\begin{enumerate}

% Source map: 2025 LR Section 2, Question 01 — point at issue
\item Daria says, ``The city should convert the vacant lot into a garden because nearby apartments lack outdoor space.'' Noel replies, ``Residents already have access to the river park, so a garden on that lot is not needed.''

Daria and Noel disagree about whether
\begin{enumerate}
\item[(A)] residents use the river park.
\item[(B)] the lot is large enough for a garden.
\item[(C)] apartment residents value outdoor space.
\item[(D)] a garden on the vacant lot is needed.
\item[(E)] the river park contains gardens.
\end{enumerate}

% Source map: 2025 LR Section 2, Question 02 — strengthen
\item A decommissioned server contains login records encrypted with a protocol unfamiliar to the investigators. They hypothesize that someone accessed the machine after it was decommissioned.

Which fact most strengthens the hypothesis?
\begin{enumerate}
\item[(A)] The server stored several kinds of files.
\item[(B)] The login records were written using software unavailable before the server was decommissioned.
\item[(C)] Other servers remained in service longer.
\item[(D)] Some records on the machine are unencrypted.
\item[(E)] The server has a metal case.
\end{enumerate}

% Source map: 2025 LR Section 2, Question 03 — resolve discrepancy
\item A factory added a second inspection to its production line, yet customers reported defective units at about the same rate as before.

Which fact best resolves the apparent discrepancy?
\begin{enumerate}
\item[(A)] The second inspection takes less than a minute.
\item[(B)] Customers submit defect reports online.
\item[(C)] The factory produces more than one model.
\item[(D)] At the same time, the factory changed to a cheaper component that created defects at roughly the rate the new inspection prevented them.
\item[(E)] Inspectors wear protective glasses.
\end{enumerate}

% Source map: 2025 LR Section 2, Question 04 — principle support
\item A proposed digital catalog meets every formatting standard, but tests show that its indexing scheme prevents intended users from finding the collection. A librarian argues that completing it unchanged would not count as a successful cataloging project.

Which principle most supports the planner's reasoning?
\begin{enumerate}
\item[(A)] Technical compliance is not sufficient for success when a project's central practical purpose is not achieved.
\item[(B)] Digital projects should always use the newest file format.
\item[(C)] Ease of use is the only relevant catalog value.
\item[(D)] User tests are always perfectly accurate.
\item[(E)] A catalog should contain fewer than one thousand records.
\end{enumerate}

% Source map: 2025 LR Section 2, Question 05 — main conclusion
\item The university should prioritize an open-source captioning tool. It handles the languages used in most courses, instructors can correct its transcripts, and the university can adapt it without waiting for a vendor. Its interface is currently plain, but that does not prevent students from reading captions.

What is the main conclusion?
\begin{enumerate}
\item[(A)] Captioning tools should support every language.
\item[(B)] Instructors can correct transcripts.
\item[(C)] Vendor software is never adaptable.
\item[(D)] A plain interface is always accessible.
\item[(E)] The university should prioritize the open-source captioning tool.
\end{enumerate}

% Source map: 2025 LR Section 2, Question 06 — criticism of reasoning
\item Inez calls a recent serialized audio program the city's first podcast. Omar objects that an older radio series was later offered as downloadable episodes. Inez replies that only a program created for on-demand subscription counts as a podcast.

Inez's reply criticizes Omar's reasoning on the ground that Omar
\begin{enumerate}
\item[(A)] assumes that the radio recordings still exist.
\item[(B)] uses ``podcast'' for a program later distributed as downloads, while Inez reserves the term for programs designed for on-demand subscription.
\item[(C)] provides no date for the recent audio program.
\item[(D)] overlooks the possibility that radio series can have episodes.
\item[(E)] assumes that every city has a radio station.
\end{enumerate}

% Source map: 2025 LR Section 2, Question 07 — flaw
\item Kira uses an encrypted messaging application because she values privacy. Mateo uses the same application. Therefore, Mateo must value privacy.

The reasoning is flawed because it
\begin{enumerate}
\item[(A)] assumes that anyone who values privacy avoids all public communication.
\item[(B)] assumes that anyone using an application owns a phone.
\item[(C)] overlooks the possibility that Mateo uses the application for a different reason.
\item[(D)] treats privacy as a necessary condition for communication.
\item[(E)] infers that Mateo never uses another application.
\end{enumerate}

% Source map: 2025 LR Section 2, Question 08 — principle support
\item A grant reviewer has collaborated with one applicant several times and owns shares in that applicant's company. Even if the reviewer believes she can remain neutral, she should withdraw from scoring that application.

Which principle most strongly supports the argument?
\begin{enumerate}
\item[(A)] Reviewers should disclose conflicts only after decisions are announced.
\item[(B)] A decision maker should withdraw when substantial personal ties create a reasonable concern about impartiality.
\item[(C)] Owning shares always prevents a person from making any rational judgment.
\item[(D)] A collaborator should automatically receive a grant.
\item[(E)] Grant panels should exclude all applicants with commercial interests.
\end{enumerate}

% Source map: 2025 LR Section 2, Question 09 — necessary assumption
\item In a design contest, some students stayed late to repair another team's machine after their own score could no longer improve. Researchers conclude that these students sometimes acted to benefit others rather than themselves.

Which assumption is required?
\begin{enumerate}
\item[(A)] The repaired team later won the contest.
\item[(B)] Students never learn by watching repairs.
\item[(C)] Helping produced no hidden self-directed benefit, such as eligibility for a separate teamwork prize.
\item[(D)] Every student in the contest stayed late.
\item[(E)] Design students cooperate more often than engineering students.
\end{enumerate}

% Source map: 2025 LR Section 2, Question 10 — flaw
\item A technician advises writers to back up important drafts immediately because their shared drive is unstable. A critic replies that the advice is wrong because repairing the drive would be better than making backups.

The critic's reasoning is questionable because it
\begin{enumerate}
\item[(A)] attacks the technician's character.
\item[(B)] treats immediate advice for people facing a present risk as though it were a competing proposal for eliminating that risk.
\item[(C)] assumes that drafts can be copied.
\item[(D)] infers that all shared drives are unstable.
\item[(E)] concludes that drive repairs never succeed.
\end{enumerate}

% Source map: 2025 LR Section 2, Question 11 — principle application
\item Principle: A researcher should not promise a collaborator access that can be provided only by violating a data-use agreement.

Which argument is most justified by the principle?
\begin{enumerate}
\item[(A)] A researcher should not promise a collaborator an embargoed dataset when the governing agreement forbids sharing it before publication.
\item[(B)] A researcher should avoid promising a meeting because meetings require scheduling.
\item[(C)] A librarian should not recommend a book that some readers dislike.
\item[(D)] A collaborator may receive any data that a majority of the team has seen.
\item[(E)] A technician should publish every preliminary measurement because delay is inconvenient.
\end{enumerate}

% Source map: 2025 LR Section 2, Question 12 — evaluate the argument
\item A school reports that students using a new tutoring program improved more quickly than students using the existing program. It concludes that the new program is more effective.

Which question would be most useful in evaluating the conclusion?
\begin{enumerate}
\item[(A)] Are tutoring programs difficult to schedule?
\item[(B)] Did the two student groups have comparable skill levels before tutoring began?
\item[(C)] Does the school offer activities other than tutoring?
\item[(D)] What colors appear in the program's logo?
\item[(E)] Do teachers prefer tutoring to independent study?
\end{enumerate}

% Source map: 2025 LR Section 2, Question 13 — resolve discrepancy
\item A theater's weekday ticket price increased, but weekday attendance also increased, even though no new productions opened.

Which fact best resolves the apparent discrepancy?
\begin{enumerate}
\item[(A)] Weekend attendance remained steady.
\item[(B)] The theater sells snacks.
\item[(C)] At the same time, the city ended free evening parking on weekends but kept it on weekdays.
\item[(D)] Some patrons buy tickets online.
\item[(E)] The theater's largest auditorium has a balcony.
\end{enumerate}

% Source map: 2025 LR Section 2, Question 14 — weaken
\item Engineers applied a new protective coating to steel bridge samples, but after six months the coated samples showed no less surface rust than untreated samples. They conclude that the coating provides no protective benefit.

Which fact most weakens the conclusion?
\begin{enumerate}
\item[(A)] Steel is used in structures other than bridges.
\item[(B)] The coating sharply reduced microscopic cracks that typically cause serious damage only after several years.
\item[(C)] The coating was developed in a university laboratory.
\item[(D)] Some bridges are painted bright colors.
\item[(E)] The test facility also stores aluminum samples.
\end{enumerate}

% Source map: 2025 LR Section 2, Question 15 — point at issue
\item Priya says, ``The laboratory should retain raw sensor data for ten years because future researchers may need to verify the published analysis.'' Tomas replies, ``Ten-year storage would consume funds needed to keep the sensors operating.''

Priya and Tomas are most clearly committed to disagreeing about whether
\begin{enumerate}
\item[(A)] the laboratory collects sensor data.
\item[(B)] future researchers will use every file.
\item[(C)] data storage costs money.
\item[(D)] published analyses should be verifiable.
\item[(E)] the laboratory should retain the raw data for ten years.
\end{enumerate}

% Source map: 2025 LR Section 2, Question 16 — principle match
\item A scientist says: We should distinguish lacking evidence that a comet contains ice from having evidence that it contains no ice. Failure to prove a claim is not ordinarily proof of its negation.

Which case most closely illustrates the principle?
\begin{enumerate}
\item[(A)] A mechanic finds evidence that a battery is charged and concludes that it is not empty.
\item[(B)] A biologist finds no reliable survey showing whether otters inhabit an inlet and therefore suspends judgment rather than declaring that no otters inhabit it.
\item[(C)] A gardener sees no rain and waters the plants.
\item[(D)] A voter rejects a proposal after reading evidence against it.
\item[(E)] A judge treats two independent witnesses as stronger evidence than one.
\end{enumerate}

% Source map: 2025 LR Section 2, Question 17 — most strongly supported
\item Standardized test scores can reveal whether students have mastered certain skills, but they do not measure every educational goal. Scores also reflect differences in preparation outside school. Thus, a score is evidence about learning, though not a complete measure of it.

Which statement is most strongly supported?
\begin{enumerate}
\item[(A)] Standardized tests reveal nothing about learning.
\item[(B)] Educational goals cannot be assessed at all.
\item[(C)] Outside preparation determines every score.
\item[(D)] Test scores may provide some information about learning without measuring it completely.
\item[(E)] Schools should consider only test scores.
\end{enumerate}

% Source map: 2025 LR Section 2, Question 18 — role in argument
\item Some historians argue that every major political transition can ultimately be explained only by individual leaders' decisions. But knowing every legal move available to each chess piece would not by itself explain why a coordinated strategy succeeded. Likewise, facts about individuals may be insufficient to explain organized social patterns. Therefore, the proposed individual-only explanation is not established.

What role does the example about grammar and a novel play?
\begin{enumerate}
\item[(A)] It is evidence that leaders have no influence on history.
\item[(B)] It gives an analogy supporting the possibility that facts about individual parts do not fully explain coordinated patterns.
\item[(C)] It defines political transitions as board games.
\item[(D)] It is the conclusion the argument seeks to refute.
\item[(E)] It shows that the rules of chess are false.
\end{enumerate}

% Source map: 2025 LR Section 2, Question 19 — sufficient assumption
\item If the data center loses either of its two backup servers, it will lack required redundancy. Server R will shut down unless its cooling unit is replaced. The cooling unit will not be replaced. Therefore, the center will lack required redundancy.

Which assumption makes the conclusion follow logically?
\begin{enumerate}
\item[(A)] Server R is one of the data center's two backup servers.
\item[(B)] The center currently has more redundancy than required.
\item[(C)] No server other than R has a cooling unit.
\item[(D)] A cooling unit can be replaced only once.
\item[(E)] Both backup servers were installed at the same time.
\end{enumerate}

% Source map: 2025 LR Section 2, Question 20 — parallel flaw
\item Advertisement: Anyone who runs a successful bakery needs reliable ovens. You have reliable ovens, so your bakery will be successful.

Which argument displays the same flaw?
\begin{enumerate}
\item[(A)] Successful gardens receive enough water. This garden receives enough water, so it will be successful.
\item[(B)] Any copper wire conducts electricity. This wire is copper, so it conducts electricity.
\item[(C)] A valid ticket is required to enter. Kai has no valid ticket, so Kai cannot enter.
\item[(D)] If the alarm sounds, the guard calls. The alarm sounds, so the guard calls.
\item[(E)] A sealed jar keeps air out. This jar is not sealed, so air may enter.
\end{enumerate}

% Source map: 2025 LR Section 2, Question 21 — flaw
\item A columnist argues: On days when ice-cream sales are high, swimming-pool attendance is also high. Therefore, buying ice cream must cause people to visit swimming pools.

The reasoning is flawed because it
\begin{enumerate}
\item[(A)] fails to consider that another factor, such as hot weather, might increase both ice-cream sales and pool attendance.
\item[(B)] assumes that pools permit visitors to bring food.
\item[(C)] relies on an appeal to expert opinion.
\item[(D)] generalizes from a single pool visitor.
\item[(E)] concludes that ice cream is never healthy.
\end{enumerate}

% Source map: 2025 LR Section 2, Question 22 — necessary assumption
\item The archive's new search tool will reduce the time historians spend locating documents because it returns results faster than the old tool.

Which assumption is required by the argument?
\begin{enumerate}
\item[(A)] Every historian searches for the same documents.
\item[(B)] The new tool costs less than the old tool.
\item[(C)] Historians spend more time reading than searching.
\item[(D)] The new tool's results are sufficiently relevant that historians will not lose the saved time sorting unusable results.
\item[(E)] The old tool will be deleted immediately.
\end{enumerate}

% Source map: 2025 LR Section 2, Question 23 — strengthen
\item After a museum added a clearly marked self-guided route, the number of visitors reaching its upper gallery increased. The director concludes that the route caused the increase by making navigation easier.

Which fact most strengthens the conclusion?
\begin{enumerate}
\item[(A)] Attendance at museums nationwide also rose that month.
\item[(B)] The museum simultaneously moved its most popular exhibit upstairs.
\item[(C)] Upper-gallery visitors frequently reported that the marked route helped them find the stairs.
\item[(D)] Marked routes are used in many buildings.
\item[(E)] The old signs were easier to clean.
\end{enumerate}

% Source map: 2025 LR Section 2, Question 24 — flaw
\item Researchers find that firms with mandatory remote-work policies retain fewer new employees than firms without such policies. They conclude that ending any mandatory remote-work policy will restore retention to its previous level.

The reasoning is most vulnerable because it
\begin{enumerate}
\item[(A)] assumes that employees never work remotely by choice.
\item[(B)] ignores the possibility that lasting workplace changes or other causes could prevent retention from returning after the policy ends.
\item[(C)] treats retention as identical to employee satisfaction.
\item[(D)] bases a prediction on comparative evidence rather than a definition.
\item[(E)] presumes that every firm hires new employees.
\end{enumerate}

% Source map: 2025 LR Section 2, Question 25 — cannot be true
\item A symposium has three speakers: Gia, Arun, and Bo. Exactly two speak in the afternoon session. If Gia speaks in the afternoon, Bo does not. If Arun does not speak in the afternoon, Gia does.

Which situation cannot be true?
\begin{enumerate}
\item[(A)] Gia and Arun speak in the afternoon.
\item[(B)] Arun and Bo speak in the afternoon.
\item[(C)] Arun speaks in the afternoon.
\item[(D)] Gia speaks and Bo does not speak in the afternoon.
\item[(E)] Gia and Bo speak in the afternoon.
\end{enumerate}

\end{enumerate}

\subsection*{Answer Key and Concise Explanations}

\subsubsection*{Practice Set 1}
\begin{enumerate}
\item \textbf{B.} The institution's stated merit-and-safety criteria, not popularity, should control.
\item \textbf{D.} Without the assumption that managers follow the conditional order, their decision would not reveal the stated belief.
\item \textbf{C.} Mina says the article is needed; Rowan says it is not.
\item \textbf{B.} Failure to prove leakage does not prove that leakage never occurs with a tightened lid.
\item \textbf{A.} If red is from 1990 or earlier while blue is later than 1990, blue is newer.
\item \textbf{D.} Machine alignment gives the holes a current function independent of hanging folders.
\item \textbf{C.} A new rail connection can increase ridership despite shorter operating hours.
\item \textbf{B.} Absorption supports reduced flooding, which supports funding.
\item \textbf{E.} Reduced traffic offers a competing cause for the decline in collisions.
\item \textbf{A.} Both arguments illicitly move from some members of a broad class to some members of a subclass.
\item \textbf{D.} More frogs can exist while fewer fall inside the survey area.
\item \textbf{B.} The comparison helps isolate ventilation from a broader change affecting nearby studios.
\item \textbf{C.} The other statements support this claim about useful diagrams.
\item \textbf{A.} Publishing evidence makes the reasons inspectable by those asked to accept the decision.
\item \textbf{D.} Sato accepts and Bell rejects that proposed tradeoff.
\item \textbf{B.} The analogy shows that a forecast may account for the process it predicts.
\item \textbf{C.} Closure plus inadequate stored oak entails missing the deadline.
\item \textbf{E.} Both are valid applications of modus tollens.
\item \textbf{A.} The trial contrast supports a stronger cold-weather constraint under rapid discharge.
\item \textbf{D.} This is the claim supported by the two benefits of temporary work.
\item \textbf{C.} The association does not establish causation or exclude another explanation.
\item \textbf{B.} Accreditation directly entails employing a veterinarian.
\item \textbf{A.} Algebra is necessary for club membership here, not sufficient for it.
\item \textbf{E.} Replacing the strap preserves and improves the borrowed camera's functionality.
\item \textbf{D.} A majority in a limited course sample cannot establish a universal claim about all evening students.
\end{enumerate}

\subsubsection*{Practice Set 2}
\begin{enumerate}
\item \textbf{D.} Daria supports and Noel denies the need for the proposed garden.
\item \textbf{B.} Software unavailable before decommissioning places the login record afterward.
\item \textbf{D.} A newly introduced source of defects can offset defects prevented by the inspection.
\item \textbf{A.} The principle connects the catalog's central practical purpose to project success.
\item \textbf{E.} The remaining statements are reasons or qualifications supporting the recommendation.
\item \textbf{B.} Their disagreement turns on whether later distribution or original design fixes the category.
\item \textbf{C.} A shared behavior can have different causes.
\item \textbf{B.} Substantial ties create the reasonable impartiality concern described.
\item \textbf{C.} If helping carried a hidden personal reward, the other-directed interpretation would not follow.
\item \textbf{B.} Immediate personal advice and system-wide prevention answer different questions.
\item \textbf{A.} The promised access would require violating the stated data-use agreement.
\item \textbf{B.} Baseline differences could explain the observed improvement.
\item \textbf{C.} The parking change can shift patrons from weekends to weekdays despite the price rise.
\item \textbf{B.} Preventing slow-developing microscopic damage is a benefit not captured by six-month surface rust.
\item \textbf{E.} Priya recommends ten-year retention while Tomas gives a reason against it.
\item \textbf{B.} The biologist avoids treating lack of confirming evidence as proof of absence.
\item \textbf{D.} The passage presents scores as informative but incomplete.
\item \textbf{B.} The analogy supports limits on explaining coordinated patterns only through facts about individuals.
\item \textbf{A.} It connects R's shutdown to the stated loss-of-redundancy condition.
\item \textbf{A.} Both affirm the consequent by treating a necessary condition as sufficient.
\item \textbf{A.} The correlation does not establish the proposed direction or rule out a common cause.
\item \textbf{D.} Faster return alone saves no total time if the results require extra sorting.
\item \textbf{C.} Visitor reports connect the marked route to reaching the upper gallery.
\item \textbf{B.} Ending the policy need not reverse every other cause or persistent workplace change.
\item \textbf{E.} Exactly two speakers are scheduled, but Gia's presence requires Bo's absence.
\end{enumerate}

